\section{Kết luận và mở rộng}

\subsection{Kết quả đạt được}
Bài tập lớn đã hoàn thành đầy đủ các mục tiêu đề ra:
\begin{enumerate}
    \item Thiết lập và giải thành công hệ phương trình vi phân bậc 2 mô tả chuyển động ném xiên có lực cản tuyến tính, sử dụng thư viện SymPy cho lời giải symbolic chính xác.
    
    \item Tính toán và phân tích chi tiết 5 khía cạnh:
    \begin{itemize}
        \item Quỹ đạo $y(x)$ cho 5 góc ném khác nhau ($15°$, $30°$, $45°$, $60°$, $75°$)
        \item So sánh định lượng quỹ đạo có và không có lực cản
        \item Biến thiên vận tốc $v_x(t)$, $v_y(t)$ theo thời gian
        \item Biến thiên năng lượng (động năng, thế năng, cơ năng, năng lượng mất)
        \item Ảnh hưởng của hệ số cản $h$ đến quỹ đạo
    \end{itemize}
    
    \item Xây dựng ứng dụng web tương tác (Streamlit + Plotly) cho phép nhập tham số bằng slider và quan sát kết quả biến đổi theo thời gian thực.
\end{enumerate}

\subsection{Nhận xét vật lý quan trọng}
Từ kết quả mô phỏng, có thể rút ra các kết luận vật lý quan trọng sau:

\begin{enumerate}
    \item Lực cản tuyến tính $\vec{F}_c = -h\vec{v}$ làm quỹ đạo \textbf{mất tính đối xứng parabol}: nhánh đi xuống dốc hơn nhánh đi lên do vận tốc ngang suy giảm theo hàm mũ.
    
    \item \textbf{Góc ném tối ưu} (cho tầm xa cực đại) dịch chuyển về phía nhỏ hơn $45°$. Trong điều kiện mô phỏng ($h = 0.5$ kg/s), góc tối ưu $\approx 30°$.
    
    \item Vận tốc ngang giảm theo $e^{-ht/m}$, vận tốc đứng tiến tới \textbf{vận tốc giới hạn} $v_\infty = -mg/h$, phù hợp với lý thuyết vật lý.
    
    \item \textbf{Cơ năng không bảo toàn} do lực cản là lực không bảo toàn, luôn sinh công âm. Năng lượng cơ học chuyển hóa thành nhiệt năng.
    
    \item Hệ số cản $h$ ảnh hưởng \textbf{phi tuyến} đến quỹ đạo -- tầm xa giảm rất nhanh khi $h$ tăng.
\end{enumerate}

\subsection{So sánh lý thuyết và mô phỏng}
Nghiệm giải tích từ SymPy khớp hoàn toàn với lý thuyết vật lý:
\begin{itemize}
    \item Khi $h \to 0$: nghiệm hội tụ về dạng parabol cổ điển (phương trình \ref{eq:ideal}).
    \item Khi $t \to \infty$: $v_x \to 0$, $v_y \to -mg/h$ (vận tốc giới hạn).
    \item Năng lượng: $E_0 = KE + PE + W_{cản}$ luôn thỏa mãn tại mọi thời điểm, xác nhận tính nhất quán của mô hình.
\end{itemize}

\subsection{Hướng mở rộng}
Một số hướng phát triển tiềm năng nếu muốn nâng cao độ chính xác của mô hình:
\begin{itemize}
    \item \textbf{Lực cản bậc hai:} $\vec{F}_{cản} = -c|\vec{v}|\vec{v}$, phù hợp hơn với vận tốc cao. Mô hình này không có nghiệm giải tích đơn giản, cần sử dụng phương pháp Runge-Kutta (\texttt{scipy.integrate.solve\_ivp}).
    \item \textbf{Hiệu ứng gió:} Thêm thành phần lực ngoài theo phương ngang.
    \item \textbf{Mô phỏng 3D:} Mở rộng sang không gian 3 chiều.
    \item \textbf{Tối ưu hóa góc ném:} Sử dụng \texttt{scipy.optimize} để tìm góc cho tầm xa cực đại ứng với mỗi bộ tham số.
\end{itemize}
